% GENERATED FILE. Edit canonical lessons, then run npm run generate:books.
% canonical-source-hash: fnv1a64:2856875171e0dc28
% canonical-lessons: AR-C538-tatriz, AR-C538-tarraza, AR-C538-hiyaka, AR-C538-haka-knit, AR-C538-khata

\chapter{Needle and thread}
\label{ch:ar-needle-and-thread}
\hlchaptermodality{538}

\begin{chapteropening}
\textbf{By the end of this chapter:} \emph{I can talk about embroidery and knitting, and say that someone embroidered, knitted or sewed.}

Complete the last lesson of chapter 538: I can talk about embroidery and knitting, and say that someone embroidered, knitted or sewed.
\end{chapteropening}

\section[\texorpdfstring{\ar{تطريز} (taṭrīz)}{تطريز (taṭrīz)}]{\ar{تطريز} (taṭrīz) — embroidery}

\label{lesson:AR-C538-tatriz}



\begin{quote}
Before the new one: say the Arabic for pottery, then the Arabic for a crossword.
\end{quote}

\subsection*{You'll want to know: \ar{تطريز}}
\textbf{\ar{تطريز}} — \emph{taṭrīz} — \textquotedblleft{}embroidery\textquotedblright{}.

\begin{quote}
\ar{تعلمت} \ar{جدتي} \ar{التطريز} \ar{في} \ar{طفولتها}.
\end{quote}

\begin{quote}
\emph{taʿallamat jaddatī at-taṭrīz fī ṭufūlatihā.}
\end{quote}

\begin{quote}
My grandmother learned embroidery in her childhood.
\end{quote}

\begin{cousinweb}
\textbf{\ar{تطريز}} (\emph{taṭrīz}) is built on the native Arabic root \emph{ṭ-r-z}, \textquotedblleft{}to embroider, to figure a cloth\textquotedblright{} (Lane). Taṭrīz is the verbal noun of Form II.
\end{cousinweb}

\begin{grammarlens}[title={What you've built}]
One more word for describing what you lived and felt.
\end{grammarlens}

\subsection*{Your turn}
\begin{itemize}
\raggedright
  \item \emph{Say it:} \emph{taṭrīz}
  \item \emph{Say it:} \emph{taṭrīz}, once more
  \item \emph{Say it:} \emph{taʿallamat jaddatī at-taṭrīz fī ṭufūlatihā.} — My grandmother learned embroidery in her childhood.
  \item \emph{Recall:} say the Arabic for pottery, then the Arabic for a crossword, then say \emph{taṭrīz} again
\end{itemize}

\begin{tcolorbox}[breakable,colback=teal!4,colframe=teal!35!black,title={Before you move on}]
What does \ar{تطريز} mean? (\textquotedblleft{}Embroidery\textquotedblright{}.) Say it once more.
\end{tcolorbox}
\section[\texorpdfstring{\ar{طرز} (ṭarraza)}{طرز (ṭarraza)}]{\ar{طرز} (ṭarraza) — to embroider}

\label{lesson:AR-C538-tarraza}



\begin{quote}
Before the new one: say the Arabic for a crossword, then the Arabic for embroidery.
\end{quote}

\subsection*{You'll want to know: \ar{طرز}}
\textbf{\ar{طرز}} — \emph{ṭarraza} — \textquotedblleft{}to embroider\textquotedblright{}.

\begin{quote}
\ar{طرزت} \ar{أمي} \ar{وردة} \ar{على} \ar{الوسادة}.
\end{quote}

\begin{quote}
\emph{ṭarrazat ummī warda ʿalā al-wisāda.}
\end{quote}

\begin{quote}
My mother embroidered a rose on the pillow.
\end{quote}

Present: \ar{يطرز}, \emph{yuṭarriz}.

\begin{cousinweb}
\textbf{\ar{طرز}} (\emph{ṭarraza}) is built on the native Arabic root \emph{ṭ-r-z}, \textquotedblleft{}to embroider, to figure a cloth\textquotedblright{} (Lane). Form II ṭarraza is the usual verb.
\end{cousinweb}

\begin{grammarlens}[title={What you've built}]
One more word for describing what you lived and felt.
\end{grammarlens}

\subsection*{Your turn}
\begin{itemize}
\raggedright
  \item \emph{Say it:} \emph{ṭarraza}
  \item \emph{Say it:} \emph{ṭarraza}, once more
  \item \emph{Say it:} \emph{ṭarrazat ummī warda ʿalā al-wisāda.} — My mother embroidered a rose on the pillow.
  \item \emph{Recall:} say the Arabic for a crossword, then the Arabic for embroidery, then say \emph{ṭarraza} again
\end{itemize}

\begin{tcolorbox}[breakable,colback=teal!4,colframe=teal!35!black,title={Before you move on}]
What does \ar{طرز} mean? (\textquotedblleft{}To embroider\textquotedblright{}.) Say it once more.
\end{tcolorbox}
\section[\texorpdfstring{\ar{حياكة} (ḥiyāka)}{حياكة (ḥiyāka)}]{\ar{حياكة} (ḥiyāka) — knitting, weaving}

\label{lesson:AR-C538-hiyaka}



\begin{quote}
Before the new one: say the Arabic for embroidery, then the Arabic for to embroider.
\end{quote}

\subsection*{You'll want to know: \ar{حياكة}}
\textbf{\ar{حياكة}} — \emph{ḥiyāka} — \textquotedblleft{}knitting, weaving\textquotedblright{}.

\begin{quote}
\ar{تحب} \ar{خالتي} \ar{الحياكة} \ar{في} \ar{الشتاء}.
\end{quote}

\begin{quote}
\emph{tuḥibbu khālatī al-ḥiyāka fī ash-shitāʾ.}
\end{quote}

\begin{quote}
My aunt loves knitting in the winter.
\end{quote}

\begin{cousinweb}
\textbf{\ar{حياكة}} (\emph{ḥiyāka}) is built on the native Arabic root \emph{ḥ-w-k}, \textquotedblleft{}to weave\textquotedblright{} (Lane). Lane gives ḥiyāka as one of the nouns of weaving.
\end{cousinweb}

\begin{grammarlens}[title={What you've built}]
One more word for describing what you lived and felt.
\end{grammarlens}

\subsection*{Your turn}
\begin{itemize}
\raggedright
  \item \emph{Say it:} \emph{ḥiyāka}
  \item \emph{Say it:} \emph{ḥiyāka}, once more
  \item \emph{Say it:} \emph{tuḥibbu khālatī al-ḥiyāka fī ash-shitāʾ.} — My aunt loves knitting in the winter.
  \item \emph{Recall:} say the Arabic for embroidery, then the Arabic for to embroider, then say \emph{ḥiyāka} again
\end{itemize}

\begin{tcolorbox}[breakable,colback=teal!4,colframe=teal!35!black,title={Before you move on}]
What does \ar{حياكة} mean? (\textquotedblleft{}Knitting, weaving\textquotedblright{}.) Say it once more.
\end{tcolorbox}
\section[\texorpdfstring{\ar{حاك} (ḥāka)}{حاك (ḥāka)}]{\ar{حاك} (ḥāka) — to knit, to weave}

\label{lesson:AR-C538-haka-knit}



\begin{quote}
Before the new one: say the Arabic for to embroider, then the Arabic for knitting.
\end{quote}

\subsection*{You'll want to know: \ar{حاك}}
\textbf{\ar{حاك}} — \emph{ḥāka} — \textquotedblleft{}to knit, to weave\textquotedblright{}.

\begin{quote}
\ar{حاكت} \ar{جدتي} \ar{وشاحا} \ar{من} \ar{الصوف}.
\end{quote}

\begin{quote}
\emph{ḥākat jaddatī wishāḥan min aṣ-ṣūf.}
\end{quote}

\begin{quote}
My grandmother knitted a woollen scarf.
\end{quote}

Present: \ar{يحوك}, \emph{yaḥūk}.

\begin{cousinweb}
\textbf{\ar{حاك}} (\emph{ḥāka}) is built on the native Arabic root \emph{ḥ-w-k}, \textquotedblleft{}to weave\textquotedblright{} (Lane). Lane gives yaḥūku as its present.
\end{cousinweb}

\begin{grammarlens}[title={What you've built}]
One more word for describing what you lived and felt.
\end{grammarlens}

\subsection*{Your turn}
\begin{itemize}
\raggedright
  \item \emph{Say it:} \emph{ḥāka}
  \item \emph{Say it:} \emph{ḥāka}, once more
  \item \emph{Say it:} \emph{ḥākat jaddatī wishāḥan min aṣ-ṣūf.} — My grandmother knitted a woollen scarf.
  \item \emph{Recall:} say the Arabic for to embroider, then the Arabic for knitting, then say \emph{ḥāka} again
\end{itemize}

\begin{tcolorbox}[breakable,colback=teal!4,colframe=teal!35!black,title={Before you move on}]
What does \ar{حاك} mean? (\textquotedblleft{}To knit, to weave\textquotedblright{}.) Say it once more.
\end{tcolorbox}
\section[\texorpdfstring{\ar{خاط} (khāṭa)}{خاط (khāṭa)}]{\ar{خاط} (khāṭa) — to sew}

\label{lesson:AR-C538-khata}



\begin{quote}
Before the new one: say the Arabic for knitting, then the Arabic for to knit.
\end{quote}

\subsection*{You'll want to know: \ar{خاط}}
\textbf{\ar{خاط}} — \emph{khāṭa} — \textquotedblleft{}to sew\textquotedblright{}.

\begin{quote}
\ar{خاط} \ar{الخياط} \ar{زرا} \ar{على} \ar{قميصي}.
\end{quote}

\begin{quote}
\emph{khāṭa al-khayyāṭ zirran ʿalā qamīṣī.}
\end{quote}

\begin{quote}
The tailor sewed a button on my shirt.
\end{quote}

Present: \ar{يخيط}, \emph{yakhīṭ}.

\begin{cousinweb}
\textbf{\ar{خاط}} (\emph{khāṭa}) is built on the native Arabic root \emph{kh-y-ṭ}, \textquotedblleft{}to sew\textquotedblright{} (Lane). Khayṭ, a thread, and khayyāṭ, a tailor, which you know, share the root.
\end{cousinweb}

\begin{grammarlens}[title={What you've built}]
One more word for describing what you lived and felt.
\end{grammarlens}

\subsection*{Your turn}
\begin{itemize}
\raggedright
  \item \emph{Say it:} \emph{khāṭa}
  \item \emph{Say it:} \emph{khāṭa}, once more
  \item \emph{Say it:} \emph{khāṭa al-khayyāṭ zirran ʿalā qamīṣī.} — The tailor sewed a button on my shirt.
  \item \emph{Recall:} say the Arabic for knitting, then the Arabic for to knit, then say \emph{khāṭa} again
\end{itemize}

\begin{tcolorbox}[breakable,colback=teal!4,colframe=teal!35!black,title={Before you move on}]
What does \ar{خاط} mean? (\textquotedblleft{}To sew\textquotedblright{}.) Say it once more.
\end{tcolorbox}
